\documentclass[11pt]{article}
\usepackage[a4paper,margin=1in]{geometry}
\usepackage{kotex}
\usepackage{amsmath}
\usepackage{booktabs}
\usepackage{array}
\usepackage{hyperref}
\usepackage{enumitem}
\usepackage{longtable}

\hypersetup{
  colorlinks=true,
  linkcolor=blue,
  urlcolor=blue,
  pdftitle={MuJoCo-ArduSub Depth Hold 문제 분석 및 수정 보고서},
  pdfauthor={Codex}
}

\title{MuJoCo JSON SITL 환경에서\\ArduSub Depth Hold 문제 분석 및 수정 보고서}
\author{local workspace report}
\date{2026-04-05}

\begin{document}
\maketitle

\section{문서 목적}

이 문서는 MuJoCo + ArduSub + JSON SITL 구조에서 발생한 \texttt{ALT\_HOLD}
/ Depth Hold 이상 동작을 어떻게 추적했고, 최종적으로 무엇을 수정해서
해결했는지를 한국어로 정리한 문서이다.

초기 증상은 다음과 같았다.

\begin{itemize}[leftmargin=1.5em]
\item \texttt{ALT\_HOLD} 전환 후 기체가 계속 하강하려고 함
\item 상승 명령을 줘도 충분히 올라오지 않거나, 하강 쪽 출력이 우세함
\item 바닥이나 목표 깊이 근처에서도 제어 출력이 정지하지 않음
\end{itemize}

\section{최종 원인 요약}

최종 원인은 ``깊이 센서를 못 읽어서''가 아니라, ArduSub가 부팅할 때
\texttt{AHRS\_EKF\_TYPE=3} 상태로 올라가면서 내부 수직 상태 추정값이
시뮬레이터 진실값과 분리된 것이었다.

즉,

\begin{itemize}[leftmargin=1.5em]
\item \texttt{BARO Alt}와 \texttt{GPS Alt}는 비교적 정상인데
\item 제어기가 실제로 쓰는 \texttt{CTUN Alt}가 크게 발산했고
\item 그 결과 ArduSub는 실제보다 훨씬 위에 있다고 오판하여
      계속 하강 출력을 냈다.
\end{itemize}

\section{분석 순서}

\subsection{1. 조이스틱 중립 문제 여부 확인}

처음에는 단순히 조이스틱 heave/throttle 입력이 완전히 중립으로
돌아오지 않아서 \texttt{ALT\_HOLD}가 hold 상태로 못 들어가는지를 확인했다.

하지만 로그에서 \texttt{RCIN}/\texttt{RC\_CHANNELS} 값이 관련 채널에서
\texttt{1500} 근처를 유지했기 때문에, ``중립 불량''만으로 현재 증상을
설명할 수는 없었다.

\subsection{2. 브리지 프레임 정합 문제 수정}

그 다음 발견된 문제는 MuJoCo 브리지에서 깊이와 수직속도를 서로 다른
기준계로 해석하고 있던 부분이었다.

구체적으로는 다음 문제가 있었다.

\begin{itemize}[leftmargin=1.5em]
\item 깊이는 world 기준 위치로 계산됨
\item 수직속도는 body/local 성격이 섞인 \texttt{cvel} 경로를 타고 들어감
\item 결과적으로 ArduSub가 받는 \texttt{position.z}와 \texttt{velocity.z}가
      같은 물리 상태를 표현하지 않을 수 있었음
\end{itemize}

그래서 \texttt{uuv\_mujoco/v2.2/bridge/ros2\_bridge.py}에서 다음을 정리했다.

\begin{enumerate}[leftmargin=1.5em]
\item \texttt{mujoco.mj\_objectVelocity(..., flg\_local=0)} 를 사용해
      world ENU 기준 선속도를 우선 사용
\item \texttt{data.cvel}은 fallback으로만 사용하고, 그 경우에도 world로 회전해서 사용
\item depth는 \texttt{bar30\_site} 기준으로 계산
\item 수직속도도 같은 \texttt{bar30\_site} 기준 속도를 사용
\end{enumerate}

이 단계는 센서값의 기준축이 엇갈리는 문제를 제거했다.

\subsection{3. JSON sensor/servo 통신 경로 정리}

또 다른 문제는 JSON transport 경로였다.

\begin{itemize}[leftmargin=1.5em]
\item sensor packet 반환 주소가 한때 고정 포트로 바뀌어 있었음
\item 실제 ArduSub JSON 입력은 servo 송신자의 동적 endpoint를 따라감
\item servo source가 중복 해석되면 명령 출처가 섞일 수 있었음
\end{itemize}

그래서 \texttt{uuv\_mujoco/v2.2/bridge/sitl\_transport.py}에서 다음을 고쳤다.

\begin{enumerate}[leftmargin=1.5em]
\item sensor JSON은 \texttt{\_sitl\_client\_addr}로 되돌리되,
      sender를 아직 모를 때만 fallback 포트를 사용
\item servo polling을 중복되지 않게 정리
\item JSON servo와 MAVLink servo가 동시에 주 명령원이 되지 않도록 정리
\end{enumerate}

이 단계는 QGC 연결과 ArduSub sensor 수신 안정성을 회복시켰다.

\subsection{4. 바닥 근처 동작 보강}

처음 현상은 ``계속 내려간다''는 형태였기 때문에 바닥 근처 동작도 같이 보강했다.

\begin{itemize}[leftmargin=1.5em]
\item DVL altitude를 downward rangefinder 성격으로 사용
\item 바닥 접촉 근처에서는 vertical truth를 정지 쪽으로 clamp
\item \texttt{surface\_bottom\_detector.cpp}에서 SITL일 때 rangefinder proximity도
      바닥 판정 보조 조건으로 사용
\end{itemize}

이 수정은 바닥 근처 안정성에는 도움이 되었지만, 이것만으로 핵심 문제가
완전히 해결되지는 않았다.

\subsection{5. 결정적 원인 확인}

핵심은 최신 BIN 로그를 비교하면서 드러났다.

\paragraph{문제 상태 로그 \texttt{00000185.BIN}}

\begin{itemize}[leftmargin=1.5em]
\item \texttt{SURFACE\_DEPTH = -1.0}
\item \texttt{AHRS\_GPS\_USE = 1}
\item \texttt{AHRS\_EKF\_TYPE = 3}
\item \texttt{RNGFND1\_TYPE = 0}
\end{itemize}

이 상태에서

\begin{itemize}[leftmargin=1.5em]
\item \texttt{BARO Alt}는 약 \(-1.1\ \mathrm{m}\)
\item \texttt{GPS Alt}도 약 \(-1.1\ \mathrm{m}\)
\item 그런데 \texttt{CTUN Alt}는 \(18.6,\ 32.4,\ 77.8,\ 107.6,\ 298.5\ \mathrm{m}\)
      식으로 비정상적으로 증가
\item 수직 thruster 출력 \texttt{RCOU C5--C8}는 \texttt{1810--1900} 근처로 포화
\end{itemize}

즉 센서는 정상인데 제어에 쓰이는 내부 수직 상태가 틀어져 있었다.

\paragraph{수정 후 로그 \texttt{00000186.BIN}}

\begin{itemize}[leftmargin=1.5em]
\item \texttt{SURFACE\_DEPTH = -1.0}
\item \texttt{AHRS\_GPS\_USE = 0}
\item \texttt{AHRS\_EKF\_TYPE = 10}
\item \texttt{RNGFND1\_TYPE = 100}
\end{itemize}

이 상태에서는

\begin{itemize}[leftmargin=1.5em]
\item \texttt{CTUN Alt}가 약 \(-0.06 \sim -0.07\ \mathrm{m}\)로 유지
\item \texttt{BARO Alt}도 비슷한 스케일을 유지
\item \texttt{RCOU C5--C8}가 대체로 \texttt{1490--1512} 근처에 머묾
\item 이전처럼 하강 출력이 지속적으로 포화되지 않음
\end{itemize}

결론적으로 결정적 원인은 수직 상태 추정 경로였고, 그 경로는 부팅 파라미터로
제어되고 있었다.

\section{최종 수정 내용}

최종적으로 결정적인 수정은
\texttt{uuv\_mujoco/v2.2/start\_ardusub\_sitl\_mj311.sh}에서 부팅 파라미터를
강제로 주입한 것이다.

\subsection{추가한 파라미터}

\begin{center}
\begin{tabular}{>{\ttfamily}p{0.24\linewidth} p{0.10\linewidth} p{0.54\linewidth}}
\toprule
파라미터 & 값 & 적용 이유 \\
\midrule
AHRS\_EKF\_TYPE & 10 & EKF3 대신 SIM 기반 AHRS 경로를 사용해서 JSON SITL 수직 상태 발산을 막음 \\
AHRS\_GPS\_USE & 0 & AHRS 쪽 GPS 기반 위치/고도 보정을 끔 \\
RNGFND1\_TYPE & 100 & SITL 가상 rangefinder backend를 활성화함 \\
RNGFND1\_MIN & 0.05 & 최소 유효 거리 5 cm 설정 \\
RNGFND1\_MAX & 30.0 & 최대 유효 거리 30 m 설정 \\
RNGFND1\_ORIENT & 25 & 센서를 downward-facing 으로 설정 \\
SURFACE\_DEPTH & -1.0 & 표면으로 간주할 깊이 기준을 유지 \\
SURFACE\_MAX\_THR & 0.1 & 표면 근처 upward throttle을 제한 \\
\bottomrule
\end{tabular}
\end{center}

\section{파라미터 상세 의미}

\subsection{\texttt{AHRS\_EKF\_TYPE}}

이 파라미터는 ArduSub가 attitude/position 추정을 어떤 backend로 할지를
결정한다. ArduPilot 코드 주석 기준 허용값은 다음과 같다.

\begin{center}
\begin{tabular}{>{\ttfamily}p{0.14\linewidth} p{0.70\linewidth}}
\toprule
값 & 의미 \\
\midrule
0 & Disabled. EKF 기반 자세/위치 추정을 사용하지 않음 \\
2 & EKF2 사용 \\
3 & EKF3 사용 \\
10 & Sim 사용. 시뮬레이터 상태를 직접 쓰는 경로 \\
11 & ExternalAHRS 사용 \\
\bottomrule
\end{tabular}
\end{center}

이번 문제에서 중요했던 점은 다음과 같다.

\begin{itemize}[leftmargin=1.5em]
\item \texttt{3}으로 두면 EKF3가 수직 상태를 추정하는데,
      현재 JSON SITL 구성에서는 그 수직 상태가 발산했다.
\item \texttt{10}으로 두면 SIM backend를 통해 시뮬레이터가 주는 상태를
      직접 반영하므로, 이번 구성에서는 훨씬 안정적이었다.
\item \texttt{2}는 EKF2인데, 지금 문제의 본질이 ``EKF류 수직 상태가 이 구성에서
      안정적이지 않다''는 쪽이므로 근본 해결책으로 보지 않았다.
\item \texttt{11}은 외부 AHRS를 따로 쓰는 경우라 현재 구조와 맞지 않는다.
\end{itemize}

이번 수정의 핵심은 바로 이 값을 \texttt{3 -> 10}으로 바꾼 것이다.

\subsection{\texttt{AHRS\_GPS\_USE}}

ArduPilot 코드 주석 기준 의미는 다음과 같다.

\begin{center}
\begin{tabular}{>{\ttfamily}p{0.14\linewidth} p{0.70\linewidth}}
\toprule
값 & 의미 \\
\midrule
0 & GPS를 DCM navigation에 사용하지 않음. dead reckoning 중심 \\
1 & GPS를 DCM position 보정에 사용 \\
2 & GPS를 DCM position과 height 둘 다에 사용 \\
\bottomrule
\end{tabular}
\end{center}

주의할 점은 코드 주석에도 나오듯이, 이 값은 주로 DCM 기반 AHRS에 직접 영향이
있고 EKF는 별도 소스 설정을 따른다는 점이다.

이번 수정에서는 \texttt{AHRS\_EKF\_TYPE=10}과 같이 쓸 때 GPS 고도 혼합이
추가로 끼지 않도록 \texttt{0}으로 두었다.

\subsection{\texttt{RNGFND1\_TYPE}}

이 파라미터는 \#1 rangefinder backend가 어떤 종류인지를 결정한다.
값 종류는 매우 많고 센서별로 다르다. 이번 문맥에서 중요한 값은 다음이다.

\begin{center}
\begin{tabular}{>{\ttfamily}p{0.14\linewidth} p{0.70\linewidth}}
\toprule
값 & 의미 \\
\midrule
0 & 비활성화 \\
1 & 기본 아날로그 계열 rangefinder \\
10 & MAVLink rangefinder backend \\
36 & scripting backend \\
100 & SITL rangefinder backend \\
\bottomrule
\end{tabular}
\end{center}

이번 문제에서는 \texttt{0}이면 downward rangefinder가 완전히 꺼져 있어서
바닥 proximity 정보를 ArduSub가 쓸 수 없었다. 그래서 \texttt{100}으로 바꿔
MuJoCo가 보내는 가상 DVL altitude / rangefinder 경로를 ArduSub가 인식하도록 했다.

\subsection{\texttt{RNGFND1\_MIN}, \texttt{RNGFND1\_MAX}}

이 값들은 rangefinder의 유효 거리 범위를 뜻한다.

\begin{itemize}[leftmargin=1.5em]
\item \texttt{RNGFND1\_MIN=0.05}: 5 cm 미만은 유효 측정으로 보지 않음
\item \texttt{RNGFND1\_MAX=30.0}: 30 m 초과는 유효 측정으로 보지 않음
\end{itemize}

이번 구성에서는 수조 환경과 DVL altitude 범위를 고려할 때 충분히 넓은
상한값과, 바닥 접촉 근처에서 너무 작은 값이 이상치로 해석되지 않도록 하는
하한값이 필요했다.

\subsection{\texttt{RNGFND1\_ORIENT}}

이 값은 센서가 어느 방향을 보는지 뜻한다. 이번에 사용한 \texttt{25}는
\texttt{ROTATION\_PITCH\_270}이고, ArduPilot 내부에서는 downward-facing
센서로 처리된다.

즉,

\[
\texttt{RNGFND1\_ORIENT = 25}
\quad \Longrightarrow \quad
\text{아래를 보는 센서}
\]

이게 맞아야 ArduSub가 바닥 거리 측정으로 인식한다.

\subsection{\texttt{SURFACE\_DEPTH}}

ArduSub 파라미터 설명 기준으로, 이 값은 ``기체가 표면에 있다고 판단할 때
외부 압력센서가 읽는 깊이''이며 단위는 cm이다.

따라서

\begin{itemize}[leftmargin=1.5em]
\item \texttt{-1.0} 은 약 \(-1\ \mathrm{cm}\) 부근을 표면 판정 기준으로 본다는 의미
\item 너무 큰 양수나 잘못된 값이면 \texttt{ALT\_HOLD} 진입 시 표면 근처로
      오판할 수 있다
\end{itemize}

이번 사례에서는 이 값 자체가 주원인은 아니었지만, 과거 잘못된 surface
semantics가 하강을 유발했던 적이 있어서 명시적으로 \texttt{-1.0}으로
고정 유지했다.

\subsection{\texttt{SURFACE\_MAX\_THR}}

ArduSub 설명 기준으로, 기체가 표면에 가까워질수록 upward throttle을 1.0에서
이 값까지 선형으로 줄이는 제한값이다. 감쇠는 표면에서 약 1 m 아래부터 시작된다.

즉,

\begin{itemize}[leftmargin=1.5em]
\item 값이 작을수록 표면 근처 상승 thrust를 더 강하게 제한
\item 값이 크면 표면 근처에서도 upward throttle이 더 많이 허용됨
\end{itemize}

이번에는 \texttt{0.1}을 유지해서 표면 근처 상승 thrust가 과도하게 커지지 않게 했다.

\section{왜 이 수정이 효과가 있었는가}

핵심은 제어기가 사용하는 내부 높이 추정값 \(\hat{z}\)를 다시 시뮬레이터 진실값에
가깝게 만든 것이다.

\[
e_z = z^\star - \hat{z}
\]

여기서

\begin{itemize}[leftmargin=1.5em]
\item \(z^\star\): hold 목표 깊이
\item \(\hat{z}\): ArduSub 내부 추정 깊이
\end{itemize}

이다. 이전에는 \(\hat{z}\)가 크게 틀어져서 \(e_z\)가 계속 큰 음수/양수로
남았고, 그 때문에 수직 thruster가 하강 쪽으로 계속 포화됐다.

수정 후에는

\begin{itemize}[leftmargin=1.5em]
\item \texttt{CTUN Alt}가 \texttt{BARO Alt}와 비슷한 크기에 머물고
\item 수직 출력이 1500 근처로 돌아왔으며
\item Depth Hold가 더 이상 무한 하강 형태로 깨지지 않았다
\end{itemize}

\section{파일별 역할}

\begin{longtable}{>{\ttfamily}p{0.40\linewidth} p{0.52\linewidth}}
\toprule
파일 & 역할 \\
\midrule
\endfirsthead
\toprule
파일 & 역할 \\
\midrule
\endhead
uuv\_mujoco/v2.2/bridge/ros2\_bridge.py & depth와 vertical velocity 프레임 정합, DVL altitude 처리, 바닥 근처 vertical truth 안정화 \\
uuv\_mujoco/v2.2/bridge/sitl\_transport.py & JSON sensor return target 복구, sensor/servo transport 경로 정리 \\
ardupilot/ArduSub/surface\_bottom\_detector.cpp & SITL 환경에서 rangefinder proximity를 바닥 판정 보조로 사용 \\
uuv\_mujoco/v2.2/start\_ardusub\_sitl\_mj311.sh & 최종 결정적 수정. 부팅 파라미터 override를 통해 수직 상태 추정 경로를 안정화 \\
\bottomrule
\end{longtable}

\section{실무 결론}

이번 문제는 처음에는 joystick neutral, 통신 문제, depth sensor 문제,
바닥 판정 문제처럼 보였지만, 최종적으로는 ``ArduSub가 어떤 상태추정기를
쓰고 있는가''가 핵심이었다.

실제로 문제를 멈춘 조합은 다음과 같다.

\[
\texttt{AHRS\_EKF\_TYPE=10},\quad
\texttt{AHRS\_GPS\_USE=0},\quad
\texttt{RNGFND1\_TYPE=100}
\]

그리고 이 조합 위에서

\begin{itemize}[leftmargin=1.5em]
\item depth/velocity 프레임 정합
\item JSON transport 정리
\item 바닥 proximity 보강
\end{itemize}

이 함께 들어가면서 전체 Depth Hold 동작이 안정화되었다.

\section{컴파일}

이 문서는 현재 위치인 \texttt{/Users/kanghyunmin/Desktop/uuv\_sim/document/docs}
에서 다음 명령으로 컴파일할 수 있다.

\begin{verbatim}
pdflatex depth_hold_fix_report.tex
\end{verbatim}

\end{document}
